\section{Lean}\todo{section Lean}
\subsection{Was ist Lean}\todo{sub\-section Lean}
Lean\cite{10.1007/978-3-030-79876-5_37} ist beides eine funktionale Progammiersprache und ein Beweis Assistent für formale Beweise.
Das Projekt wurde 2013 von Leonardo de Moura und Sebastian Ullrich ins Leben gerufen.
Die aktuelle Version Lean 4 ist selbst in Lean geschrieben und wurde im September 2023 released.

Lean basiert stark auf seinem Typen-System
Dieses ordnet jedem Term in Lean einen Typ zu, welcher wiederum selbst ein Term ist.
Hier ein par Beispiele:
\begin{lstlisting}
"Hello World!" : String
16 : Nat
{ val := 16, isLt := by omega } : Fin 32
16 < 32 : Prop
isLt : 16 < 32
Nat : Type
Prop : Type
Fin : Nat → Type
Type : Type 1
Type 1 : Type 2
\end{lstlisting}
Ein String ist eine Zeichenkette und Nat oder auch $\mathbb{N}$ eine natürliche Zahl.
Fin n mit einer natürlichen Zahl n ist die Menge der natürlichen Zahlen, die kleiner als n sind.
Ein Element von Fin n besteht zum einen aus einer Zahl val und zum anderen aus einem Nachweis isLt, dass val kleiner als n ist.
Thesen (engl. Propositions) sind vom Typ Prop.
Dabei ist jede These selbst wider ein Typ und die Elemente dieses Typs sind die Beweise, die sie belegen.
Was Thesen von anderen Typen unterseidet, ist, dass alle ihre Elemente gleich sind.
Es macht für Lean keinen Unterschied, wie eine These bewiesen wurde, solange es mindestens einen Beweis gibt.
Alle Typen wie String, Nat, Fin 32 oder Prop haben den Typ Type.
Auch Funktionen haben Typen, nämlich den von wo nach wo sie abbilden.
Fin ist zum Beispiel vom Typ der Funktionen, die von Nat auf Type abbilden.
Da in Lean alles einen Typ hat, gilt das auch Type.
Type, was das gleiche ist wie Type 0, ist vom Typ Type 1.
Type 1 ist wiederum vom Typ Type 2 und so weiter.

Um eigene Typen beziehungsweise Terme definieren, gibt es in Lean verschiedene Möglichkeiten:
Mit def oder abbrev lassen sich vor allem Funktionen definieren, aber auch konstante Elemente eines Typs oder einfache zusammen gesetzte Typen, wie zum Beispiel Tupel.
Der Hauptunterschied dabei ist, dass aus sicht von Lean def etwas neues kreiert über das noch nichts bekannt ist, wärend abbrev eher ein Alias mit neuem Namen erzeugt.
Komplexere Typen lassen sich als structure definieren was an Klassen aus dem objektorientierten Progammieren erinnert mit beliebig vielen Atributen und Methoden.
abbrev
def
structure
inductive

Der Taktik Modus und das Laufen beim Schreiben

Für mehr Informationen über Lean und seine Geschiche können auf der Lean Webseite\cite{lean_lang} oder der Lean Sprach Referenz\cite{lean_lang_doc} nachgelesen werden.

\subsection{Anwendung}\todo{sub\-section Anwendung}
Microsoft nutzt Lean um nachweislich sichere Kryptographie zu Produzieren.
Wie in diesem Bolg-Post\cite{microsoftresearchblog} eines Microsoft-Research-Teams zu lesen, nutzt Microsoft um die Einhaltung von Standarts zu verifizieren.
Da in der Kryptographie der Code häufig sehr hardwarenah und architekturabhängig geschrieben wird
und Laufzeitoptimirungen enthält sieht er selten genau so aus, wie der Algorithmus aus dem Standart.
Um hierbei sicher zu gehen, dass die Funktionsweise immer noch die gleiche ist, reicht Testen nicht immer aus, da ein kleiner Fehler die Sicherheit des gesamten Systhems untergraben kann.
Deshalb formalizieren sie die Standarts in Lean.
Eine Implementation in Rust wird dann auf eine äquivalente Implementatierung in Lean abgebildet.
Der Vorteil gegenüber dem Testen ist, dass nun bewiesen werden kann, dass für alle möglichen Inputs, die eventuelle Vorbedingungen erfüllen, auch ein korrekter Output erzeugt wird.
Wohingegen Testen immer nur für begrenzt viele Inputs zeigen kann, dass der Output korrekt ist.

\subsection{Mathlib}\todo{sub\-section Mathlib}
\cite{mathlib2020}

\subsection{Lean und KI}\todo{sub\-section KI}

Sobalt das mathematische Problem formaliziert wurde, kann KI versuchen es zu lösen.
Oder wenn bereits ein Beweis bekannt ist, kann KI probieren einen einfacheren oder schnelleren zuu finden.

In \cite{microsoftresearchblog} wird erklärt, wie KI bereits zum Einsatz kommt und welche Aufgaben sie übernimmt.
Zum einen das Übersetzen von Standarts in Lean Spezifikationen.
Diese müssen natürlich getestet und von Menschen geprüft werden, spraren unterm Strich aber trozdem Arbeit.
Zum anderen beim automatischen Schreiben und Aktualisieren von Beweisen.
Das Schöne ist, das Lean garantiert, dass ein Beweis auch funktioniert, sofer also der Fehler nicht schon in der Formulierung der zu beweisenden Aussage liegt,
kann ein Beweis, der von einer KI geschrieben wurde, nicht fälschlich positiv Sein.
Menschen müssen nur die Beweise per Hand schreiben, bei denen die KI scheitert.
